# 4. Experiments and Results

  

本节通过三组实验验证所提摆臂式覆盖路径规划方法的效果。第一组实验将本文方法与 Fields2Cover 进行对比，用于说明主动摆臂扫掠模型相对于固定宽度条带式 CPP 的差异；第二组实验比较不同矩形候选扫描策略，用于分析矩形分解如何影响复杂空间中的覆盖质量和冗余；第三组实验改变车体速度，并固定摆臂角速度，用于揭示摆臂扫掠中速度匹配对覆盖质量和执行时间的影响。

  

三组实验均使用两张基础地图。`map_test1` 包含四个规则矩形障碍，用于观察基础绕障和覆盖能力。`parkingGraph` 包含密集分布的小型障碍，更接近被柱、孔洞和局部结构切割的施工平面。第二组实验另加入三张方向性地图：

- `scan_comb_corridors`：横向梳状通道；
- `scan_mixed_weave`：横纵交织结构；
- `scan_offset_rooms`：错位房间结构。

这三种结构用于检验行扫描、列扫描和双方向扫描的差异。

  

所有地图均为二值占据图，白色像素表示待覆盖自由区域，黑色像素表示障碍或不可作业区域；分辨率为 $0.5\ \mathrm{m/pixel}$。除速度实验外，本文方法采用统一参数：抹盘半径 $r_d=0.75\ \mathrm{m}$，摆臂长度 $L=1.20\ \mathrm{m}$，车体宽度 $0.80\ \mathrm{m}$，车体最大线速度 $2.00\ \mathrm{m/s}$，车体最大角速度 $30^\circ/\mathrm{s}$，摆臂角速度上限 $90^\circ/\mathrm{s}$。矩形候选最小边长为 $2.5\ \mathrm{m}$，贪婪选择权重为 $\alpha=0.5$、$\beta=2.0$，随机种子固定为 42。



## 4.1 Comparison with Fields2Cover

  

Fields2Cover 是适合本文的固定宽度 CPP 基线，也是当前该类问题中具有代表性的开源框架。它经过公开发表并持续开源维护，集成了多种先进的覆盖规划算法与评价目标，包含地头生成、作业条带生成、路线排序和曲率受限路径连接等完整模块；同时支持非凸区域、内部障碍物、固定作业宽度、车辆宽度、转弯半径和速度参数。从功能完整性、算法集成度和可复现性看，Fields2Cover 代表了当前固定宽度 CPP 的综合先进水平（state of the art, SOTA），因而不仅对应成熟的“固定机具宽度 + 条带覆盖 + 路径连接”规划范式，也构成具有竞争力的工程基线。本文方法与其对比，目标不是比较两个方法的具体实现细节，而是突出主动摆臂扫掠模型与综合 SOTA 固定宽度条带规划框架在有障碍空间中的覆盖形态和可执行性差别。

  

实验中，Fields2Cover 的作业宽度设为 $1.50\ \mathrm{m}$，与本文半径 $0.75\ \mathrm{m}$ 抹盘的直径一致；车体宽度为 $0.80\ \mathrm{m}$，安全余量为 $0.05\ \mathrm{m}$，因此在原始自由空间上执行 $0.45\ \mathrm{m}$ 的安全内缩。两种方法均按完整轨迹计覆盖：本文方法以抹盘沿车体—摆臂轨迹的扫掠区域计算，Fields2Cover 则以固定宽度工具沿完整扫掠区域（含作业、转弯和转场）计算。为使评价范围与实际规划范围一致，Fields2Cover 的覆盖率以安全内缩后的自由区域面积为分母；在计算固定宽度工具沿轨迹形成的扫掠面积时，条带端部按平端处理，轨迹拐角按斜接处理。作业条带速度设为 $2.0\ \mathrm{m/s}$，转弯与连接路径速度保持为 $0.5\ \mathrm{m/s}$。

  

**表 2. 本文方法与 Fields2Cover 的对比结果**

```{=latex}
\begin{center}
\footnotesize
\setlength{\tabcolsep}{3pt}
\renewcommand{\arraystretch}{1.18}
\begin{tabular}{@{}p{2.1cm}p{2.0cm}rrrrp{2.7cm}@{}}
\toprule
地图 & 方法 & \shortstack{评价面积\\(m²)} & \shortstack{覆盖面积\\(m²)} & \shortstack{覆盖率\\(\%)} & \shortstack{总时长/规划时\\(s/s)} & 路径检查 \\
\midrule
map\_test1 & 本文方法 & 611.00 & 575.00 & 94.11 & 854.73/1.298 & 通过 \\
 & Fields2Cover & 533.17 & 418.00 & 78.40 & 295.50/0.075 & 越界 4.51 m² \\
parkingGraph & 本文方法 & 4294.00 & 4110.50 & 95.73 & 5844.73/17.981 & 通过 \\
 & Fields2Cover & 4046.79 & 3756.85 & 92.84 & 2883.81/6.339 & 越界 20.00 m² \\
\bottomrule
\end{tabular}
\end{center}
```

注：两种方法的评价面积不同，覆盖率不能直接作同域差值比较。Fields2Cover 的越界路线未通过路径检查，其时长和规划用时仅记录生成结果。

  

在 `map_test1` 上，使用与抹盘直径一致的 $1.50\ \mathrm{m}$ 固定作业宽度后，Fields2Cover 的生成时间和标称执行时长更低。按完整轨迹扫掠计算，其覆盖 418.00 m²（78.40%），仍低于本文方法的 575.00 m²（94.11%）；但连接段存在 4.51 m² 的车体扫掠越界。因此，295.50 s 的时长和 0.075 s 的规划时间只能作为无效路径的生成结果记录，不能作为可执行方案的性能成绩。该结果体现了摆臂构型的核心特点：车体不必严格沿每一条待覆盖条带行驶，抹盘可以通过摆臂运动扩展覆盖范围，从而在同一自由区域内形成更宽的时变扫掠。

  

在 `parkingGraph` 上，本文方法覆盖 4110.50 m²，覆盖率达到 95.73%；Fields2Cover 的完整轨迹覆盖 3756.85 m²，覆盖率为 92.84%。不过，$1.50\ \mathrm{m}$ 工具宽度下的连接路径越出自由空间，产生 20.00 m² 的无效车体扫掠；相应的 2883.81 s 时长和 6.339 s 规划时间同样不能作为有效基线成绩。该问题集中出现在局部障碍密集的位置：条带本身可以在内缩后的区域中生成，但不同条带之间的曲率受限连接需要额外空间，一旦缺少类似农田地头的宽阔转场区，连接段就可能穿过原始自由空间之外的区域。本文方法通过摆臂扫掠把车体运动与末端覆盖解耦，在复杂空间中保留了更大的覆盖能力和路径连续性。

  
![本文方法与 Fields2Cover 的覆盖路径和覆盖区域对比](figures/figure_2_coverage_comparison_cn.png)

**图 2. 本文方法与 Fields2Cover 的覆盖路径和覆盖区域对比。** (a)、(b) 对应 `map_test1`（图内沿用 `map_simple` 标签），(c)、(d) 对应 `parkingGraph`；左列为本文方法，右列为 Fields2Cover。蓝色表示完整轨迹扫掠区域，白色或浅色表示评价域内尚未覆盖的自由空间，黑色表示障碍或不可作业区域，红色标记表示车体碰撞或无效扫掠位置。两张地图中的 Fields2Cover 路径均未通过车体碰撞检查：`map_test1` 越界 4.51 m²，`parkingGraph` 越界 20.00 m²。因此，右列的覆盖率与时长仅展示生成结果，不能作为可执行方案成绩。

  

## 4.2 Effect of Rectangle Candidate Scan Strategy

  

上一节说明，主动摆臂模型能够在障碍环境中形成不同于固定宽度条带的覆盖效果；这一优势能否稳定发挥，还取决于矩形分解方法能否恰当地表达自由空间。若候选生成受单一扫描方向偏置，后续轨迹即使可执行，也可能经过碎片化矩形、和大量重叠矩形，造成不必要的转场。因此，本节考察候选扫描方向对“矩形集合—访问序列—完整轨迹”的传导影响。

实验比较仅行扫描（row-only）、仅列扫描（col-only）和先行后列扫描（row-then-col）。除候选生成方向外，机器人参数、最小矩形边长、贪婪选择权重、随机种子、连通性修复、访问排序和局部覆盖生成均保持一致。评价指标包括贪婪覆盖率、最终分区覆盖率、分区重叠率、全轨迹覆盖率和总时长，用以区分分区表达质量与实际执行性能。实验流程为：从二值图生成候选矩形并进行贪婪选择，完成补洞和连通修复后，再以相同的轨迹生成和圆盘扫掠评估方法计算全轨迹覆盖率与任务时间。

  

**表 3. 三种候选扫描策略在五张地图上的结果**

  

```{=latex}
\begingroup
\footnotesize
\setlength{\tabcolsep}{2.5pt}
\renewcommand{\arraystretch}{1.10}
\begin{longtable}{@{}p{2.9cm}p{2.2cm}>{\raggedleft\arraybackslash}p{1.6cm}>{\raggedleft\arraybackslash}p{1.75cm}>{\raggedleft\arraybackslash}p{1.5cm}>{\raggedleft\arraybackslash}p{1.6cm}>{\raggedleft\arraybackslash}p{1.6cm}@{}}
\toprule
地图 & 策略 & \shortstack{贪婪\\覆盖率} & \shortstack{最终分区\\覆盖率} & \shortstack{分区\\重叠率} & \shortstack{全轨迹\\覆盖率} & \shortstack{总时长\\(s)} \\
\midrule
\endfirsthead
\toprule
地图 & 策略 & \shortstack{贪婪\\覆盖率} & \shortstack{最终分区\\覆盖率} & \shortstack{分区\\重叠率} & \shortstack{全轨迹\\覆盖率} & \shortstack{总时长\\(s)} \\
\midrule
\endhead
map\_test1 & row-only & 72.01\% & 100.00\% & 34.21\% & 93.04\% & 696.89 \\
 & col-only & 72.01\% & 100.00\% & 34.21\% & 93.04\% & 696.89 \\
 & row-then-col & 100.00\% & 100.00\% & 0.00\% & 94.11\% & 854.73 \\
\midrule
parkingGraph & row-only & 87.32\% & 100.00\% & 75.34\% & 95.73\% & 5844.73 \\
 & col-only & 87.32\% & 100.00\% & 75.34\% & 95.90\% & 5896.03 \\
 & row-then-col & 87.32\% & 100.00\% & 75.34\% & 95.73\% & 5844.73 \\
\midrule
梳状通道 & row-only & 87.18\% & 100.00\% & 34.46\% & 93.82\% & 2065.09 \\
 & col-only & 86.53\% & 100.00\% & 31.46\% & 93.94\% & 2032.55 \\
 & row-then-col & 96.22\% & 100.00\% & 27.13\% & 94.01\% & 1986.42 \\
\midrule
交织区域 & row-only & 84.23\% & 100.00\% & 47.89\% & 95.26\% & 1532.65 \\
 & col-only & 83.17\% & 100.00\% & 52.59\% & 95.49\% & 1575.63 \\
 & row-then-col & 96.76\% & 100.00\% & 40.74\% & 95.70\% & 1522.77 \\
\midrule
错位房间 & row-only & 82.95\% & 100.00\% & 38.78\% & 91.97\% & 1496.73 \\
 & col-only & 81.81\% & 100.00\% & 36.74\% & 91.51\% & 1443.31 \\
 & row-then-col & 99.14\% & 100.00\% & 24.46\% & 92.06\% & 1374.38 \\
\midrule
五图平均 & row-only & — & — & — & 93.96\% & 2327.22 \\
 & col-only & — & — & — & 93.98\% & 2328.88 \\
 & row-then-col & — & — & — & 94.32\% & 2316.60 \\
\bottomrule
\end{longtable}
\endgroup
```


  

首先考察两张一般地图，以判断三种候选策略在非定向测试场景中的基本表现。`map_test1` 中，row-then-col 以 94.11% 获得最高全轨迹覆盖率，并将分区重叠率降至 0，但其总时长为 854.73 s，高于 row-only 和 col-only 的 696.89 s。因此，在这一规则障碍场景中，双方向策略以更长的执行时间换取了更高的覆盖率。`parkingGraph` 中，row-then-col 与 row-only 选中了相同的矩形，所得结果也完全一致；col-only 的覆盖率略高，为 95.90%，而 row-only 和 row-then-col 以 5844.73 s 并列取得最短总时长。两张一般地图说明，row-then-col 并不会在每个场景和每项指标上都优于单方向策略：当自由空间具有明显主方向或行向候选已足以表达其结构时，row-only 或 col-only 可能获得更短时间或略高覆盖率。

在上述一般场景结果的基础上，三张方向性地图进一步用于有针对性地检验扫描方向偏置，而不用于证明双方向策略在任意地图上均为最优。上述三张方向性地图分别强化了横向通道、横纵交织和错位连接等结构，使单一扫描方向留下的残余区域更容易被观察。在这三张图上，row-then-col 均取得最高贪婪覆盖率和最低分区重叠率，说明列向候选能够补充行向候选未充分表达的区域，而非仅靠增加矩形数量提高覆盖。该改进也传导到执行层面：row-then-col 在三张图上均取得最高全轨迹覆盖率和最短总时长，结果依次为 94.01% / 1986.42 s、95.70% / 1522.77 s 和 92.06% / 1374.38 s，三图平均为 93.93% / 1627.86 s。由此可见，当地图中确实同时存在横向和纵向残余结构时，行、列互补的候选集能够缓解单一扫描的方向偏置，并同时减少冗余访问、提高完整轨迹的有效扫掠。

从五图平均看，row-then-col 的全轨迹覆盖率为 94.32%，高于 row-only 的 93.96% 和 col-only 的 93.98%；平均总时长为 2316.60 s，也略低于 row-only 的 2327.22 s 和 col-only 的 2328.88 s。因而，row-then-col 更恰当的定位是稳健的默认候选生成策略：当地图同时含有横向和纵向残余结构时，它更可能保持较高覆盖率并避免明显的时间损失；对于主方向单一或结构规则的区域，仍可根据覆盖率与工期的具体权衡选择单向策略。图 3 以 `scan_mixed_weave` 为例，直观展示了这一方向互补关系。

![`scan_mixed_weave` 地图上的候选扫描策略对比](figures/figure_3_segmentation_order_en.png)

**图 3. `scan_mixed_weave` 地图上的候选扫描策略对比。** (a) 原始地图；(b) row-only 分区；(c) col-only 分区；(d) row-then-col 分区。该地图同时包含横向和纵向狭长自由空间，单方向候选会保留不同方向的残余结构；row-then-col 则通过顺序补充候选，形成更完整且重叠更低的分区表达。

  

## 4.3 Effect of Vehicle Forward Speed–Swing-Arm Angular Velocity Coupling on Coverage Performance

  

前两个实验主要回答“路径和分区是否合理”，但摆臂式覆盖还存在一个固定宽度 CPP 中不突出的执行问题：即使几何路线相同，不同车体速度也会产生不同的抹盘扫掠密度。车体前进越快，摆臂在相邻摆扫之间来不及覆盖的前进距离越大；车体前进越慢，覆盖更充分但任务时间显著增加。因此，若只报告一条几何路径，而不分析速度与摆臂角速度的匹配关系，就无法说明该路径在实际执行中对应的覆盖质量和时间代价。

本实验用于验证第 3.5.3 节中的速度耦合分析：在摆臂角速度固定时，车体速度升高应当缩短总执行时间，但也会增加相邻摆扫之间的间距并降低覆盖率。实验的目标是在相同矩形集合和相同访问顺序下隔离“车速”这一因素，观察覆盖质量和时间效率如何随速度变化。

具体地，实验固定摆臂角速度为 $\omega_a=90^\circ/\mathrm{s}$，在 `map_test1` 和 `parkingGraph` 两张地图上分别执行同一规划结果，并将车体速度从 0.5 m/s 增加到 3.0 m/s。每个速度设置下，矩形分区、矩形访问顺序和局部作业结构保持不变，仅改变车体沿轨迹执行时的速度参数。为便于把线速度与摆臂角速度放在同一尺度上比较，定义速度-摆臂角速度比为

  

$$
\rho=\frac{v}{\omega_a},
$$

  

其单位为 m/°，表示摆臂每转动 $1^\circ$ 时车体前进的距离。$\rho$ 越大，相邻摆扫之间的前进间距越大，覆盖空隙越容易出现。实验记录每组速度下的覆盖率、总时长和单位覆盖面积耗时，用于观察“覆盖质量”和“时间效率”之间的权衡。

  

```{=latex}
\newpage
```

**表 4. 不同车体速度下的全轨迹覆盖率与执行时间**

  

```{=latex}
\begin{center}
\footnotesize
\setlength{\tabcolsep}{3pt}
\renewcommand{\arraystretch}{1.12}
\begin{tabular}{@{}p{2.8cm}>{\raggedleft\arraybackslash}p{1.65cm}>{\raggedleft\arraybackslash}p{2.05cm}>{\raggedleft\arraybackslash}p{2.2cm}>{\raggedleft\arraybackslash}p{2.1cm}>{\raggedleft\arraybackslash}p{2.8cm}@{}}
\toprule
地图 & \shortstack{车速\\(m/s)} & \shortstack{$\rho$\\(m/°)} & \shortstack{全轨迹\\覆盖率} & \shortstack{总时长\\(s)} & \shortstack{单位面积耗时\\(s/m²)} \\
\midrule
map\_test1 & 0.5 & 0.00556 & 96.44\% & 1680.45 & 2.852 \\
 & 1.0 & 0.01111 & 95.46\% & 1095.53 & 1.878 \\
 & 1.5 & 0.01667 & 88.83\% & 928.81 & 1.711 \\
 & 2.0 & 0.02222 & 88.75\% & 829.74 & 1.530 \\
 & 2.5 & 0.02778 & 86.25\% & 772.26 & 1.465 \\
 & 3.0 & 0.03333 & 85.27\% & 735.12 & 1.411 \\
\midrule
parkingGraph & 0.5 & 0.00556 & 99.15\% & 14525.14 & 3.412 \\
 & 1.0 & 0.01111 & 96.65\% & 8346.21 & 2.011 \\
 & 1.5 & 0.01667 & 93.74\% & 6214.22 & 1.544 \\
 & 2.0 & 0.02222 & 92.33\% & 5130.64 & 1.294 \\
 & 2.5 & 0.02778 & 91.03\% & 4526.86 & 1.158 \\
 & 3.0 & 0.03333 & 90.32\% & 4120.11 & 1.062 \\
\bottomrule
\end{tabular}
\end{center}
```


  

表 4 的结果表明，车速提高仍会带来明确的覆盖--时间权衡，但新版全轨迹覆盖率显示该权衡比原先判断的更为温和。两张地图的总时长均随车速单调下降：`map_test1` 从 1680.45 s 降至 735.12 s，减少 56.25%；`parkingGraph` 从 14525.14 s 降至 4120.11 s，减少 71.63%。相应地，$\rho$ 从 0.00556 m/° 增至 0.03333 m/° 时，全轨迹覆盖率分别由 96.44% 降至 85.27% 和由 99.15% 降至 90.32%，仅分别下降 11.17 和 8.83 个百分点。也就是说，较大的摆扫间距会削弱覆盖连续性，但在所测速度范围内并未使任一地图的全轨迹覆盖率降至不可用水平。

  

单位全轨迹覆盖面积耗时进一步揭示了这一取舍。该指标在两张地图上均随车速增加而持续下降：`map_test1` 从 2.852 降至 1.411 s/m²，`parkingGraph` 从 3.412 降至 1.062 s/m²，分别改善 50.53% 和 68.88%。因此，3.0 m/s 同时给出最短总时长和最低单位覆盖面积耗时，但分别以 85.27% 和 90.32% 的覆盖率为代价；0.5 m/s 则给出两图的最高覆盖率。若任务要求两图均保持不低于 95% 的全轨迹覆盖率，1.0 m/s 是本组测试中可满足该约束的最高车速，且相较 0.5 m/s 已将总时长分别缩短 34.81% 和 42.54%。因此，车速不应作为轨迹生成后的附属参数，而应依据最低覆盖率约束和工期目标，与摆臂角速度共同确定：覆盖质量优先时选择较小的 $\rho$，时间效率优先且允许约 9--11 个百分点覆盖损失时可选择较大的 $\rho$。

  
![`map_test1` 上的速度--摆臂角速度比实验结果](figures/figure_5a_vehicle_speed_map_simple_en.png)

![`parkingGraph` 上的速度--摆臂角速度比实验结果](figures/figure_5b_vehicle_speed_parkingGraph_en.png)

**图 4. 速度--摆臂角速度比对全轨迹覆盖率和执行时长的影响。** (a) `map_test1`（图内沿用 `map_simple` 标签）；(b) `parkingGraph`。横轴为 $\rho=v/\omega_a$，上方辅助坐标给出对应车速。蓝色实线表示全轨迹覆盖率，红色实线表示总任务时长；两条曲线均基于表 4 的新版汇总数据绘制。

  

## 4.4 Discussion

  

三组实验共同指向同一个核心结论：本文方法的关键并不在于采用了另一种区域分解方式，而在于将车体移动与末端抹盘摆扫的协同关系贯穿地图分解、轨迹生成和实际覆盖评价，使覆盖结果由抹盘在连续执行过程中的真实扫掠决定。与 Fields2Cover 的对比说明，固定宽度条带在简单场景中高效，但在障碍密集场景中覆盖能力和连接可执行性都会受到限制；本文方法通过摆臂末端的时变扫掠扩大了有效覆盖空间。候选扫描实验说明，矩形分解必须服务于地图结构，双方向候选能够减少单一扫描方向造成的表达偏置。速度实验进一步说明，摆臂式 CPP 必须同时考虑路径几何和执行速度，车速与摆臂角速度的匹配会直接决定覆盖质量。

  

这些结果也说明，摆臂机器人并不是在所有指标上无条件优于固定宽度平台。对于规则开阔区域，固定条带方法可以以更短路径和更低时间完成可接受覆盖；摆臂的价值主要体现在固定车体难以贴近的边界、障碍密集区域和需要减少车体动作的湿混凝土表面。因此，实际施工中可以根据场景结构选择平台和参数：开阔大面区域适合固定宽度高效扫线，复杂边角和多障碍区域更适合摆臂扫掠，二者也可以在同一楼层任务中分区协同。

  

当前方法仍有若干限制。首先，矩形候选和访问序列采用启发式生成，能够提供可行且可解释的路径，但不保证全局最优；当地图中存在大量细长通道时，矩形数量和借道访问仍可能增加。其次，本文以二维静态地图为输入，尚未显式处理施工过程中的动态人员、材料堆放、定位漂移和混凝土状态变化。再次，速度分析主要考虑车体线速度与摆臂角速度的覆盖密度关系，尚未纳入更完整的底盘动力学、抹盘压力、表面质量反馈和控制误差。总体而言，本文实验表明，面向摆臂式抹平机器人时，显式建模末端扫掠、区域分解和速度耦合是区别于传统固定宽度 CPP 的关键；这些限制则为进一步走向真实工地部署留下了明确方向。

  

